% Einheit 1 von 8 – Ereignisse als Mengen (bank/zufallsexperimente-und-pfadregeln/e1.jsonl, zusammenbau v0.1)
%% TODO Zeitmarke Einheit 1: Marken-Zeile nicht lesbar
%% TODO Zweigzeile Teil 1: was hier gelernt wird (Satz in Schülersprache aus dem Einheitstitel) – Einheit 1
\einheitenkopf[e1]{Einheit 1 von 8 · Ereignisse als Mengen}
\zweigzeile{Abitur GK}
\setcounter{aufgabe}{11}

%% TODO Titel: Ich-kann-Satz fehlt in der Bank (Platzhalter: Kettenname) – Nr. 12
%% TODO Anweisung: ein Satz über der ersten Teilaufgabe fehlt in der Bank – Nr. 12
\begin{aufgabe}{Ereignisse als Mengen}
%% TODO \rechenplatz{n}: Schrittzahl des Grundfalls fehlt in der Bank (nur bei fester Schreibform, 2.3 b)
\begin{teile}
\teil Zwei Würfe mit einem Würfel: „Erst eine 2, dann eine 6.“ Ergebnis oder Ereignis?\\ \kreuz{einzelnes Ergebnis}\\ \kreuz{Ereignis aus mehreren Ergebnissen}
\teil Eine Münze wird dreimal geworfen (K = Kopf, Z = Zahl). Ergebnismenge als Tupel? Laplace-Experiment?
\teil Eine Münze wird dreimal geworfen. A: „mindestens zweimal Kopf“, B: „der erste Wurf ist Zahl“. Ergebnisse von $A \cap B$ und von „weder A noch B“?
\teil Ein Würfel und ein Glücksrad mit den Zahlen 1, 2 und 3 werden benutzt, die Zahlen addiert. Alle Ergebnisse (Würfel; Rad) mit der Summe 5?
\teil Unter den Schülern einer Schule: S = „fährt mit dem Rad“, H = „trägt einen Helm“. Deute $P(S \cap H) = 0{,}18$ und $P_S(H) = 0{,}45$.
\teil Im Tierheim: H = „ist ein Hund“, J = „ist jünger als zwei Jahre“; $P(H) = 0{,}45$, $P(J) = 0{,}3$, $P(H \cap J) = 0{,}1$. Beschreibe $\overline{H \cup J}$ und weise nach: $P(\overline{H \cup J}) = 0{,}35$.
\teil Urlaubspläne: M = „ans Meer“, B = „in die Berge“; $P(M) = 0{,}45$, $P(B) = 0{,}35$, $P(M \cap B) = 0{,}15$. P(entweder Meer oder Berge)? P = \leerfeld
\teil Pendler: R = „fährt mit dem Rad“, U = „nutzt die U-Bahn“. Deute $P(R) + P(U) - 2 \cdot P(R \cap U)$ im Sachzusammenhang. \hfill (Abitur 2023 LK)
\end{teile}
\end{aufgabe}

%% TODO Titel: Ich-kann-Satz fehlt in der Bank (Platzhalter: Kettenname) – Nr. 13
%% TODO Anweisung: ein Satz über der ersten Teilaufgabe fehlt in der Bank – Nr. 13
\begin{aufgabe}{Ereignisse als Mengen – Fehler finden, Begründen, Anwendung, Darstellungswechsel}
\begin{teile}
\teil $P(A) = 0{,}4$, $P(B) = 0{,}25$, $P(A \cap B) = 0{,}1$. Gesucht: P(entweder A oder B). Tom rechnet $0{,}4 + 0{,}25 - 0{,}1 = 0{,}55$. Finde den Fehler.
\teil Begründe: Beim Additionssatz wird $P(A \cap B)$ einmal abgezogen, bei „genau eines von beiden“ zweimal.
\teil Mobilität in einer Stadt: 62\,\% der Haushalte haben ein Fahrrad, 48\,\% ein Auto, 25\,\% beides. Für ein Carsharing-Angebot zählen Haushalte ohne beides. Wie viel Prozent sind das? \leerfeld[\%]
\teil R = „fährt Rad“, B = „fährt Bus“. Schreibe mit Symbolen: „Die Person fährt Rad, aber nicht Bus.“
\end{teile}
\end{aufgabe}
